\documentclass[11pt]{article}
\usepackage[T1]{fontenc}
\usepackage[utf8]{inputenc}
\usepackage{lmodern,microtype}
\usepackage[a4paper,margin=28mm]{geometry}
\usepackage{amsmath,amssymb,amsthm,mathtools,booktabs,longtable,array}
\usepackage[dvipsnames]{xcolor}
\usepackage{xr-hyper}
\usepackage[colorlinks=true,linkcolor=MidnightBlue,citecolor=MidnightBlue,urlcolor=MidnightBlue]{hyperref}
\newtheorem{theorem}{Theorem}[section]
\newtheorem{proposition}[theorem]{Proposition}
\newtheorem{lemma}[theorem]{Lemma}
\newtheorem{corollary}[theorem]{Corollary}
\theoremstyle{definition}\newtheorem{definition}[theorem]{Definition}
\theoremstyle{remark}\newtheorem{remark}[theorem]{Remark}
\newtheorem{question}[theorem]{Question}
\newtheorem{conjecture}[theorem]{Conjecture}
\DeclareMathOperator{\FPdim}{FPdim}
\DeclareMathOperator{\Irr}{Irr}
\DeclareMathOperator{\Tr}{Tr}
\DeclareMathOperator{\Rep}{Rep}
\let\Vec\relax
\DeclareMathOperator{\Vec}{Vec}
\newcommand{\C}{\mathcal C}
\newcommand{\Z}{\mathcal Z}
\newcommand{\Q}{\mathbb Q}
\newcommand{\N}{\mathbb Z_{\geq0}}
\newcommand{\R}{\mathbb R}
\newcommand{\CC}{\mathbb C}
\newcommand{\ZZ}{\mathbb Z}
\newcommand{\id}{\operatorname{id}}
\allowdisplaybreaks[2]
\setlength{\emergencystretch}{2em}

\hypersetup{pdftitle={Verification guide: rank-five rings through multiplicity twenty}}
\title{Verification guide\\Rank-five rings through multiplicity twenty}
\date{23 September 2026}
\begin{document}\maketitle
\section{What this release certifies}
The pinned \texttt{fusion-ring-census} repository \cite{Census} records
$40999$ pairwise nonisomorphic rank-five fusion rings through
multiplicity twenty. The checkers validate the release tensors against
the ring axioms. Exactly
seventeen match the constructed unitary models. The other $40982$
have verified general exclusion certificates. In particular, the two
formerly unresolved tensors at local indices $5073$ and $5076$ at
exact multiplicity eighteen now have general certificates: a new
inert-prime parity theorem forces any categorification to be spherical,
and their existing exact spherical-character certificates then exclude
it. Thus general, pivotal, spherical, pseudo-unitary and unitary
categorification is classified through twenty. The old signed-search
timeouts remain discovery diagnostics; they are not used as certificates.

\section{Environment and complete verification commands}
Use Python 3.11 or later, SymPy and NumPy, and a C++17 compiler supporting
OpenMP and \texttt{\_\_int128}. The recorded run used Python 3.13.5;
\texttt{provenance/environment.json} records exact versions and the
compiler. SciPy is required only for certificate discovery. The normal
verifier does not call linear programming or accept any floating-point
certificate. Install the declared Python packages and run, from the
archive root,
\begin{verbatim}
python -m pip install -r requirements.txt
python verify.py --jobs 2
python software/verify_extension.py --jobs 2
python -O verify.py --jobs 2
python -O software/verify_extension.py --jobs 2
\end{verbatim}
The commands check every certificate, not just a sample or the final
frontier. The optimized invocation confirms that no proof condition depends
on a removable \texttt{assert}. A nonzero exit or an exception is failure,
not partial certification. Each command writes a summary JSON and a
per-certificate result file in \texttt{verification/}. The supplied logs
are measurements, not substitutes for rerunning the checker.

The optional \texttt{--skip-regressions} flag omits only the secondary
controls; it still checks the full census, all exclusions, all models
and both inert-parity/spherical composite certificates. It is not the command used
for the archived full runs. Use \texttt{--output PATH} to preserve the
archived summary when running a local copy. Local timing may vary with
hardware and process contention.

\section{Census input and archival labels}
The finite classification uses the pinned census repository
\cite{Census}; its rank-five counts continue through multiplicity twenty.
The pinned rank-five tensor and count files are bundled under
\texttt{data/pinned\_repository/rank5/}. The local file
\texttt{data/classification\_tensors.json} retains the basis choices
used by the exclusion certificates. The verifier checks both the pinned
file hash and an exact isomorphism-class crosswalk before checking any
categorification certificate. The complete tensors and fingerprints
address those certificates; there is no second census computation.

Every local row stores a full tensor \texttt{N}, multiplicity \texttt{m},
historical or appended label \texttt{j}, and local ordinal
\texttt{local\_index}. Existing Vercleyen--Slingerland matches
retain the first historical source index and basis. At multiplicity
eighteen no source rows were present, so the two index fields agree.
Always select a tensor by its full multiplication coefficients if a
citation or lookup crosses data releases.

The first full verification command checks the crosswalk automatically.
The extension command checks the pinned repository order directly for
every exact-multiplicity-nineteen and -twenty tensor and every general
exclusion certificate. The files \texttt{data/decisions\_all.json} and
\texttt{.csv} collect the $40999$ decisions; \texttt{repository\_index}
refers only to direct pinned-source records, whereas \texttt{j} and
\texttt{local\_index} in the historical segment retain their
archival meanings.
To rerun just this comparison using the bundled copy of the pinned
repository file, run
\begin{verbatim}
python software/compare_pinned_census.py \
    --report verification/census_crosswalk.json
\end{verbatim}
The comparison minimizes every complete tensor under all $24$
unit-preserving basis permutations. The included report records equal
sets for all multiplicities through eighteen, together with SHA-256
hashes of both input files. Raw labeled-tensor fingerprints are not
invariant under a basis permutation.

\section{The exclusion certificate contract}
The dictionary \texttt{certificates/exclusions.json} is keyed by the
SHA-256 fingerprint of the complete stored tensor. A fingerprint is an
address only: the verifier recomputes the tensor's hypotheses and the
certificate. It cannot establish nonexistence by itself. No exclusion
criterion accepts \texttt{m}, \texttt{j} or a previous status as input.

The arithmetic leaves contain a factor/divisibility witness, finite-prime
cyclotomic obstruction or nonintegral fusion-subring index. Positive
branches verify weak integrality or the two-FP-values reduction, then
central induction. Quadratic-parity certificates recover $A_i$ from the
square rows, verify all FP multiplication equations, squarefreeness,
conductor congruence, the Casimir polynomial, its mod-two zero multiplicity,
the square-row kernel and irreducibility. Inert-prime parity certificates
add an exhaustive four-element-field character check and central
divisibility filter; they recompute the exact spherical obstruction.
No fixed radicand or list of labels is assumed by either criterion.

A spherical certificate exhausts the irreducible character components
of the commutative rational algebra. A general signed certificate first
checks that spherical certificate, then regenerates every associative
signed tensor in the positive-second-indicator normalization and every
remaining dimension-character component. Timeouts, omitted branches,
incorrect field factors or missing low matrices cause verification to
fail. The parity theorems replace this search where their hypotheses
hold; they do not assert pivotality for arbitrary rank-five rings.

For low induction the checker recomputes every bounded column of required
dimension and every row-sum completion. Native Gram pruning discards
only a negative entry or principal minor. High columns satisfy the full
rank-one Gram inequality and the algebraic dimension tests. In a
singular remainder the kernel is imposed before a nonsingular principal
block is used. The joint trace box includes an all-zero type for all
larger orders. This avoids falsely bounding a category's exponent by
its moment height.

For every low matrix and permitted second indicator the finite-indicator
certificate supplies an integer vector $y$. After reconstructing the
complete integer system $Ax=b$, the verifier checks
$A^{\mathsf T}y\ge0$ and $b^{\mathsf T}y<0$. Rational and quadratic
fields use the same code path; all available rational indicator bounds
through height six are included. The cyclic-four criterion uses the
same relative traces after the proved mixed-support and central-involution
reductions. Its balanced Galois branches instead check an explicitly
nonintegral fusion coefficient. Every support and indicator branch is
accounted for.

\section{Applying the criterion to another tensor}
The discovery driver does not look up the decision ledger. It either
matches a constructed positive model or applies the same tensor criteria
and exactly verifies the returned certificate. For example,
\begin{verbatim}
python software/discover.py --multiplicity 18 --index 5073 \
    --height 6 --seconds 120 \
    --output verification/inert_example.jsonl
\end{verbatim}
The index in this command is the local ordinal stored in
\texttt{local\_index}. An arbitrary
external rank-five tensor can instead be supplied as a JSON list of
five $5\times5$ integer matrices:
\begin{verbatim}
python software/discover.py --tensor new_tensor.json \
    --height 6 --seconds 120 --output verification/new.jsonl
\end{verbatim}
A JSON object with an \texttt{N} or \texttt{tensor} field is also accepted.
For the displayed example, inert-prime parity and the spherical-character
certificate apply automatically from the tensor; the historical index
has no role in the criterion. A separate independent residue-field
enumeration is available as
\texttt{python software/audit\_inert\_residues.py}.
For other inputs the total per-tensor allowance is separate from the signed subprocess
allowance, which is half of it. An unsupported field, absent separating
ray, exception or timeout produces \texttt{unresolved}, not
\texttt{excluded}. Raising the budget or height is permitted; a new
successful result must still pass the exact checker.

To continue beyond twenty, take the next verified tensor set from the
repository census \cite{Census}, construct row records by the local
canonical convention, and pass that JSON through \texttt{--census}. Its mathematical criteria
do not need new index-specific functions. \texttt{export\_tables.py}
regenerates the present release's manuscript tables and CSV ledger only
after a successful full verification. A future release must explicitly
update its endpoint assertions and rerun all certificates.

\section{Exactness, controls and scope of trust}
All proof arithmetic is integral or rational, including factorization
and real-root interval comparisons. NumPy is used for bounded integer
tensor operations, not eigenvalue evidence: with rank five and bound
thirty-two each multiplication sum is at most $5\cdot32^2$.
The signed generator computes a bound for each polynomial before using
64-bit evaluations; diagnostic counters fail on overflow.
Native low-Gram inputs are bounded by $10^6$; their
four-dimensional determinant sums fit in checked 128-bit ranges.
High-Gram bounds are validated by the Python wrapper before native
execution. An overflow or failed divisibility halts the computation.

The complete runs additionally compare native Gram-vector enumeration
with arbitrary-integer Cartesian enumeration on four matrices, compare
native low joins with unpruned Python joins and exact minors on four
instances, and deliberately corrupt new
criterion certificates. The known positive $K(C_4,4)$ model has an exactly checked rational
feasible branch in the cyclic central-support system. This control invokes
no numerical linear programming. Scalar tests check the Farkas
orientation and reject floating-point rays. These are supplementary
controls, not claims of formal verification of the compiler, algebra
library or categorical proofs.

The code does not construct every unitary associator. Positive existence
is established by the cited categorical constructions; the code verifies
their fusion tensors and exact census matches. Likewise, a certificate
checker verifies applications of proved necessary conditions, not the
proofs of those conditions in a theorem prover. The mathematical proof,
software and data are separate, inspectable parts of the result.

\section{Files and rebuilding}
\texttt{README.md} gives the directory map and principal commands.
\texttt{data/decisions\_all.json} and its CSV companion record all five
existence statuses through twenty; \texttt{data/decisions.json} preserves
the historical tensor labels through eighteen.
\texttt{data/extension19\_20/full.jsonl} contains the directly indexed
repository tensors and exact certificates.
\texttt{data/frontier.json} preserves the two historically difficult tensors,
formal codegrees, FP dimensions, spherical certificates and discovery
diagnostics. Their complete general certificates are in one tensor-indexed
file. All source inputs, run measurements and version information are
recorded in \texttt{provenance/}. Obsolete case-specific scripts and
status tables are retired, not left as contradictory alternative entry
points.

The LaTeX sources and PDFs for the manuscript and this guide are in
\texttt{papers/}. Run \texttt{python build\_pdfs.py} from the
root to compile them. For current counts use the pinned repository
\cite{Census}.

\section*{AI use disclosure}
GPT-6 Astra assisted with mathematical exploration, code development, and manuscript preparation.
\begin{thebibliography}{1}
\sloppy
\bibitem{Census} \emph{Fusion-ring census},
\href{https://github.com/sebastienpalcoux/fusion-ring-census/tree/b86eb6dc6eeeafe5fc5f6906d7b074e0322e9726}{repository at commit \texttt{b86eb6dc6eeeafe5fc5f6906d7b074e0322e9726}},
\texttt{results/rank5/counts.csv} and \texttt{results/rank5/tables.txt.gz} (2026).
\end{thebibliography}
\end{document}
